% =====================================================================
\chapter{Muestreo y estimadores}
\label{tema:2}
% =====================================================================

La estadística inferencial trata los datos como valores observados de variables aleatorias cuya distribución depende de parámetros desconocidos. Este tema introduce el modelo de muestra aleatoria simple, la noción de estimador y los criterios para juzgar su calidad: sesgo, varianza, error estándar, error cuadrático medio y consistencia. Se demuestran los dos resultados centrales del muestreo: la media muestral es un estimador insesgado de la media y la varianza muestral con denominador $n-1$ lo es de la varianza. Se sigue el capítulo~6 de Wasserman~\cite{wasserman}, «Models, Statistical Inference and Learning», la §6.3 de Blitzstein y Hwang~\cite{bh}, «Sample moments», la clase~12 del curso Stats~200 de Stanford~\cite{stats200} y la lectura~10a del curso 18.05 del MIT~\cite{mit1805}.

\section{Recordatorio: esperanza y varianza}

Se recogen las propiedades de la esperanza y la varianza de variables aleatorias (discretas o continuas) que se usarán continuamente~\cite[cap.~4 y §7.3]{bh}. Todas las variables que aparecen se suponen con varianza finita.

\begin{proposicion}[Propiedades de la esperanza y la varianza]\label{prop:t2-props}
Sean $X,Y,X_1,\dots,X_n$ variables aleatorias y $a,b\in\R$.
\begin{enumerate}
  \item Linealidad: $\E[aX+bY]=a\E[X]+b\E[Y]$ y, en general, $\E\bigl[\sum_iX_i\bigr]=\sum_i\E[X_i]$.
  \item $\Var(X)=\E[(X-\E X)^{2}]=\E[X^{2}]-(\E X)^{2}$, y $\Var(aX+b)=a^{2}\Var(X)$.
  \item $\Var\bigl(\sum_iX_i\bigr)=\sum_i\Var(X_i)+2\sum_{i<j}\Cov(X_i,X_j)$.
  \item Si $X$ e $Y$ son independientes, $\E[XY]=\E[X]\E[Y]$ y, por tanto, $\Cov(X,Y)=0$. En consecuencia, si $X_1,\dots,X_n$ son independientes, $\Var\bigl(\sum_iX_i\bigr)=\sum_i\Var(X_i)$.
  \item La covarianza es bilineal,
  \[
  \Cov\Bigl(\sum_ia_iX_i,\sum_jb_jY_j\Bigr)=\sum_{i,j}a_ib_j\Cov(X_i,Y_j),
  \]
  y $\Cov(X,X)=\Var(X)$.
\end{enumerate}
\end{proposicion}

\section{Muestra aleatoria, parámetro y estadístico}

\begin{definicion}[Muestra aleatoria simple]\label{def:t2-mas}
Una \emph{muestra aleatoria simple} de tamaño $n$ de una distribución $F$ es una colección de variables aleatorias $X_1,\dots,X_n$ independientes y todas con distribución $F$; se dice que son \emph{independientes e idénticamente distribuidas} (\iid) y se escribe $X_1,\dots,X_n\sim F$ \iid\ Los datos $x_1,\dots,x_n$ son una realización (valores observados) de la muestra.
\end{definicion}

El modelo describe el muestreo con reemplazamiento de una población, o el muestreo sin reemplazamiento cuando la población es mucho mayor que la muestra, y también la repetición independiente de un experimento: lanzar una moneda, generar una respuesta de un modelo de lenguaje con una semilla distinta o elegir al azar una pregunta de un banco muy grande.

\begin{definicion}[Modelo paramétrico, parámetro y estadístico]\label{def:t2-parametro}
Un \emph{modelo paramétrico} es una familia de distribuciones $\{F_\theta:\theta\in\Theta\}$ indexada por un \emph{parámetro} $\theta$ (un número o un vector) que toma valores en el \emph{espacio paramétrico} $\Theta$. Un \emph{estadístico} es cualquier función $T=g(X_1,\dots,X_n)$ de la muestra que no depende de parámetros desconocidos; es, por tanto, una variable aleatoria, y su distribución se llama \emph{distribución muestral} de $T$.
\end{definicion}

Por ejemplo, $\{\Ber(p):p\in[0,1]\}$, $\{\Pois(\lambda):\lambda>0\}$ y $\{\Nor(\mu,\sigma^{2}):\mu\in\R,\sigma>0\}$ son modelos paramétricos. Los estadísticos más importantes son la \emph{media muestral} y la \emph{varianza muestral}
\[
\Xbar=\frac1n\sum_{i=1}^{n}X_i,\qquad S^{2}=\frac{1}{n-1}\sum_{i=1}^{n}(X_i-\Xbar)^{2},
\]
que son las cantidades del \cref{tema:1} vistas como variables aleatorias. Se reserva la mayúscula para la variable aleatoria y la minúscula ($\xbar$, $s^{2}$) para su valor observado.

\section{Estimadores y sus propiedades}

\begin{definicion}[Estimador, sesgo, error estándar y error cuadrático medio]\label{def:t2-estimador}
Un \emph{estimador} de un parámetro $\theta$ es un estadístico $\widehat\theta=\widehat\theta(X_1,\dots,X_n)$ que se usa para aproximar $\theta$; su valor con unos datos concretos es una \emph{estimación}. Se definen:
\begin{enumerate}
  \item el \emph{sesgo}: $\sesgo(\widehat\theta)=\E_\theta[\widehat\theta]-\theta$; el estimador es \emph{insesgado} si su sesgo es nulo para todo $\theta\in\Theta$;
  \item el \emph{error estándar}: $\ee(\widehat\theta)=\sqrt{\Var_\theta(\widehat\theta)}$;
  \item el \emph{error cuadrático medio}: $\ECM(\widehat\theta)=\E_\theta\bigl[(\widehat\theta-\theta)^{2}\bigr]$.
\end{enumerate}
El subíndice $\theta$ indica que las esperanzas se calculan con la distribución $F_\theta$.
\end{definicion}

El sesgo mide el error sistemático y el error estándar, la dispersión del estimador alrededor de su propia media. El error estándar suele depender de $\theta$, que es desconocido; al sustituir $\theta$ por su estimación se obtiene el \emph{error estándar estimado}, que se escribe $\widehat{\ee}$.

\begin{proposicion}[Descomposición sesgo--varianza]\label{prop:t2-ecm}
$\ECM(\widehat\theta)=\sesgo(\widehat\theta)^{2}+\Var(\widehat\theta)$.
\end{proposicion}

Un estimador ligeramente sesgado puede tener menor $\ECM$ que uno insesgado si su varianza es bastante menor (\cref{eje:t2-ecm-normal}); este compromiso entre sesgo y varianza reaparece en el aprendizaje automático (módulo~8).

\section{La media muestral}

\begin{teorema}[Media y varianza de la media muestral]\label{teo:t2-media}
Sean $X_1,\dots,X_n$ \iid\ con media $\mu$ y varianza $\sigma^{2}$. Entonces
\[
\E[\Xbar]=\mu,\qquad \Var(\Xbar)=\frac{\sigma^{2}}{n},\qquad \ee(\Xbar)=\frac{\sigma}{\sqrt n}.
\]
En particular, $\Xbar$ es un estimador insesgado de $\mu$.
\end{teorema}

El error estándar de la media decrece como $1/\sqrt n$: para dividirlo por dos hay que multiplicar el tamaño de la muestra por cuatro. Esta ley de la raíz cuadrada gobierna el coste de casi toda estimación por muestreo, desde las encuestas hasta la evaluación de modelos.

\addtocounter{definicion}{1}%

\section{La varianza muestral}

\begin{teorema}[Insesgadez de la varianza muestral]\label{teo:t2-s2}
Sean $X_1,\dots,X_n$ \iid\ con media $\mu$ y varianza $\sigma^{2}$, con $n\ge2$. Entonces $\E[S^{2}]=\sigma^{2}$.
\end{teorema}

La demostración explica el $n-1$: las desviaciones se miden respecto de $\Xbar$, que está más cerca de los datos que la verdadera media $\mu$ (es la que minimiza la suma de cuadrados), de modo que $\sum_i(X_i-\Xbar)^{2}$ es en promedio menor que $\sum_i(X_i-\mu)^{2}$, exactamente en una cantidad $\sigma^{2}$.

\begin{corolario}[Sesgo de la varianza con denominador $n$]\label{cor:t2-sesgo}
El estimador $\widehat\sigma^{2}=\frac1n\sum_i(X_i-\Xbar)^{2}=\frac{n-1}{n}S^{2}$ cumple $\E[\widehat\sigma^{2}]=\frac{n-1}{n}\sigma^{2}$, luego tiene sesgo $-\sigma^{2}/n$: subestima la varianza, aunque el sesgo tiende a cero cuando $n\to\infty$.
\end{corolario}

\begin{observacion}
Que $S^{2}$ sea insesgado para $\sigma^{2}$ no implica que $S$ lo sea para $\sigma$: en general $\E[S]<\sigma$ (\cref{eje:t2-s-sesgo}). La insesgadez no se conserva por transformaciones no lineales.
\end{observacion}

\section{Consistencia y ley de los grandes números}

Un estimador razonable debe acercarse al parámetro cuando crece el tamaño de la muestra. Se escribe $\widehat\theta_n$ para subrayar la dependencia de $n$.

\begin{definicion}[Consistencia]\label{def:t2-consistencia}
Una sucesión de variables aleatorias $Y_n$ \emph{converge en probabilidad} a una constante $c$, y se escribe $Y_n\convp c$, si para todo $\varepsilon>0$
\[
\lim_{n\to\infty}\PP\bigl(\abs{Y_n-c}>\varepsilon\bigr)=0.
\]
Un estimador $\widehat\theta_n$ es \emph{consistente} si $\widehat\theta_n\convp\theta$ para todo $\theta\in\Theta$.
\end{definicion}

\begin{lema}[Desigualdades de Markov y de Chebyshev]\label{lem:t2-chebyshev}
Si $Y\ge0$ y $a>0$, entonces $\PP(Y\ge a)\le\E[Y]/a$. En consecuencia, si $W$ tiene media $\mu_W$ y varianza finita, para todo $\varepsilon>0$
\[
\PP(\abs{W-\mu_W}\ge\varepsilon)\le\frac{\Var(W)}{\varepsilon^{2}}.
\]
\end{lema}

\begin{proposicion}[Un $\ECM$ que tiende a cero implica consistencia]\label{prop:t2-ecm-consistencia}
Si $\ECM(\widehat\theta_n)\to0$ para todo $\theta$, entonces $\widehat\theta_n$ es consistente. En particular, basta que el sesgo y la varianza tiendan a cero.
\end{proposicion}

\begin{teorema}[Ley débil de los grandes números]\label{teo:t2-lgn}
Si $X_1,X_2,\dots$ son \iid\ con media $\mu$ y varianza finita $\sigma^{2}$, entonces $\Xbar_n\convp\mu$; es decir, la media muestral es un estimador consistente de la media.
\end{teorema}

La cota de Chebyshev es válida para cualquier distribución, pero es muy conservadora; en el \cref{tema:3} el teorema central del límite dará aproximaciones mucho más precisas.

\addtocounter{definicion}{1}%

\addtocounter{definicion}{1}%

